\scanpage{183}
\mspage{495}

a u t r e vie plus intense et plus délicate, que nous sommes trop endormis ou trop pusillanimes souvent pour vivre dans la vie éveillée. Et il suffit de contempler un bébé endormi, ou seulement d'être tiré d'un sommeil profond, pour sentir que même sans rêves, le vrai bon sommeil est un t r a v a i l à sa façon: quelque chose qui nous absorbe de façon totale, à ``refaire le plein'' en somme d'une énergie qui s'était dispersée et que nous venons r e p u i s e r à sa source... C'est là, à nouveau, le ``yang dans le yin'', sans quoi le yin lui-même serait destructeur.

Des réflexions dans le même sens pourraient sûrement se développer aussi pour l'inaction é v e i l l é e , en dehors du temps de sommeil. Il n'y a qu'à observer sur pièces, de façon attentive, tel ou tel état qu'on perçoit comme ``inaction''. On se rendra compte que dans l'inaction, il y a action, fût-ce le stérile caquetage d'une pensée qui continue à tourner en rond alors qu'ell\unsure{e} a cessé de travailler. Mais à vrai dire, il est impropre d'appeler ``action'' ce mouvement-là, purement mécanique, qui se poursuit par le seul effet d'une inertie - par l'incapacité d'arrêter la machine! Et ce n'est certes pas cette agitation intérieure qui va apporter à ``l'inaction'' une harmonie yin-yang qui la rende bienfaisante. Par contre, il peut en être ainsi des activités diverses destinées à meubler ses loisirs (quand ceux-ci sont néanmoins vécus comme un état d'inaction). Mais même dans l'état de repos complet d'un état de convalescence disons, il peut y avoir action, sans quoi ce repos ou ``inaction'' devient un a v a c h i s s e m e n t , peu propice certes à une convalescence (c'est à-dire justement, au rétablissement d'un équilibre perturbé!). Par exemple, cet état de repos peut susciter une attention à son propre corps et à son entourage immédiat (qui en constitue comme une deuxième peau...), une prise de connaissance donc voire une communion, qui par elle-même a un authentique caractère d'``action''; car nul doute que d' a p p r e n d r e est bel et bien un a c t e (puisqu'il a un irrécusable e f f e t : l'apparition d'une connaissance...).

En examinant un à un les quatorze couples que j'ai inclus dans le groupe action-inaction (et sûrement on pourrait en trouver bien d'autres qui s'y insèrent naturellement), on constate que pour tous sauf peut-être un, c'est le premier terme, le terme ``masculin'', qui est investi de prestige, de ``valeur'', selon les attitudes-réflexes véhiculés par notre culture et inculqués depuis l'enfance. C'est le signe de toujours ce même déséquilibre invétéré dans notre culture,

% archive sheet number: 182
